\documentclass[11pt]{article}
\usepackage[margin=1in]{geometry}
\usepackage{amsmath,amssymb,amsfonts,amsthm,mathtools,bm}
\usepackage[T1]{fontenc}
\usepackage{microtype}
\newtheorem{theorem}{Theorem}
\newtheorem{lemma}{Lemma}
\newtheorem{cor}{Corollary}
\newtheorem{defn}{Definition}
\newtheorem{principle}{Principle}
\newtheorem{statement}{Statement}
\theoremstyle{remark}
\newtheorem{statementr}{Statement}

\title{Theory Formation and the Topology of Configuration Space}
\author{James Taylor}
\date{Working draft derived from the 2003 thesis}
\begin{document}
\maketitle

\begin{quote}\small This expository manuscript overlaps with later published work and is retained as a thesis-derived, self-published treatment with a different organization and emphasis.\end{quote}

\section{Theory formation and the influence of 
the topology of configuration space} \label{c:top}


\section{Introduction} \label{s:introi}

In this chapter, we analyze the influence of the topology of the
configuration space in formulating Bohmian theories. 
We formulate the Abelian Quantization Principle\ which states that, in addition to the
usual data such as the potential, one needs to choose
a character of the fundamental group of the configuration space 
 when forming a Bohmian theory.  
To conclude the principle, we consider three different approaches
based on: the covering space; connections; and self-adjoint extensions.
We demonstrate that each approach generates equivalent theories. 
This will set us up for the next chapter. In particular, we will 
conclude that, in the theory of
identical particles moving in 3-space, 
one has the choice between bosons and fermions. This is the usual
Bose-Fermi alternative.


The current goal  is to explain the relevance of
a space being multiply connected. The prime example to have in mind is
the Aharonov-Bohm effect. Briefly, the setup is 
a particle moving along a circle embedded in 3-space; there is also 
a magnetic field  whose support is disjoint from the circle but whose integral
over the disk does not vanish. We would like to
model this situation in terms of a Bohmian system involving just the
circle with no need to embed it into an ambient space. 

We start by considering what role the covering space might have. 
The idea is to consider wave functions
on the universal covering space such that they still define a Bohmian
motion on the base space. For the circle, the covering space is a
line; the useful wave functions are those that are $2\pi$-periodic up to a
phase. In general, we find a periodicity condition 
 characterized by a character of the fundamental
group. Another approach is to consider different connections. This is
very familiar in the context of the Aharonov-Bohm effect. The
important point is that the characters of the fundamental group are in
bijective correspondence with the flat connections; the correspondence
is that of the holonomy operator. For the circle, this is the integral
of the vector potential. The final approach is to consider
different self-adjoint extensions of the Hamiltonian. 
On the circle, if we remove a
point, then we may view the space as an interval. As is well-known,
there are a variety of possible self-adjoint Laplacians on the
interval; they are 
characterized by different boundary conditions. The extensions that
are permissible are those for which, given any other point, there is a
corresponding Hamiltonian which provides the same Bohmian evolution of
the configuration. 


The notion of multiply connected spaces giving rise to different
quantum theories is not new. A sampling of the literature is the
following. The use of a covering space was done at
least as early as 1950 as in \cite{BH50}; it was more fully done in
\cite{D72}. 
The use of vector potentials was done
 in \cite{AB59}. Path integrals on multiply connected
spaces began largely with the work of \cite{S68} and \cite{DL71}; see
\cite{Sch81} for details. 
Most of these works are dedicated to scalar-valued wave functions.


This chapter is organized as follows. 
In Section {\ref{is:circ}}, we discuss the circle and its relation to 
the Aharonov-Bohm effect. In that section, we give the principle for
Bohmian theories on manifolds for which the wave functions are complex-valued
functions over the configuration space. 
We give the
general formulation and argument for the principle in
Section {\ref{s:expo}}; the generalization refers to taking wave functions to be
sections of a given complex vector bundle over the configuration space. 
The entire discussion, centered upon deriving how
characters arise, is based upon appealing to the covering
space of the configuration space.
Section {\ref{s:fb}}\ and Section {\ref{s:sa}}\ discuss approaches that do
not use the covering space.
We  contrast the various equivalent approaches in
Section {\ref{s:comp}} and discuss in Section {\ref{is:leqv}}
 how the principle is generating
locally equivalent theories. 
We conclude with a discussion of quantization. 
In particular, quantization of 
systems on the configuration space whose potential can
only be formulated on the covering space is described in
Section {\ref{s:pers}}.  Appendix \ref{s:setten} reviews the
necessary notion of a set-indexed tensor
product. Appendix \ref{s:sat}  
deals with a technical point involving the self-adjointness of the
formal Hamiltonian. 


\section{The circle and the Aharonov-Bohm effect} \label{is:circ}

In this section, all wave functions will be complex-valued functions on the
configuration space. 

Consider the system of a single particle
constrained to travel along a circle. Thinking of the circle as
sitting in \ensuremath{\mathbb{R}^{3}}, we assume that going through the center of the 
circle is a
magnetic field. We demand that the field
 vanishes in a neighborhood around the circle, but
that its integral over a disk, with the circle as the boundary, does not
vanish. Our goal is a Bohmian theory that accurately describes this
system in which the configuration space is explicitly the circle. 

This is a simplified version of the usual Aharonov-Bohm effect. The
standard answer is that the vector potential which generates the
magnetic field should be included in the covariant derivative. This
does work as explained in Section {\ref{s:fb}}. But we want to describe
this without getting into bundles. Instead, we shall focus on the
topology of the circle and see that the possibilities naturally arising
from being multiply-connected already provides the setting for such a
model. 

We are looking for other Bohm-type theories governing an evolution of
a particle on the circle. One important constraint is that the Bohmian
velocity field needs to be well-defined on the circle. The inspiration
in this approach comes from 
noting that if we write \ensuremath{\psi}\ in the polar form
$Re^{iS/\ensuremath{\hbar}}$, then $v^{\ensuremath{\psi}} = \nabla S$. Thus, the Bohmian velocity
field is locally a gradient. From classical experiences, this
immediately brings to mind the universal covering space which, for the
circle, is just $\ensuremath{\mathbb{R}^{}}$. We can therefore ask which  wave functions on
$\ensuremath{\mathbb{R}^{}}$ could provide an evolution on the circle. The tentative
answer are those wave functions for which $S$ is periodic up to a constant,
i.e.\ $S( \theta + 2k\pi) = S(\theta) + k\beta$. This translates into a
multiplicative periodicity condition on $\ensuremath{\psi}$: 
\begin{equation} \label{ie:sper}
\ensuremath{\psi}(\theta + 2k\pi) = \gamma^k \ensuremath{\psi} (\theta).
\end{equation}
As a side note, we could say that the periodicity condition is
represented by a homomorphism between the additive group of integers
into the multiplicative group of the unit circle in the complex
plane. More abstractly, we could say that we have a unitary,
1-dimensional representation of
the covering group, and therefore of the fundamental group, for the
circle. 


Given this periodicity condition, it is clear to see that 
the associated Bohmian velocity field on \ensuremath{\mathbb{R}^{}}\ is $2\pi$-periodic. Hence,
one can project the velocity field down to the circle in order to
define a Bohm-type evolution of the particle. As
for this condition being preserved under the Schr\"odinger\ evolution on the
covering space,
linearity of Schr\"odinger's equation\ suggests it; locally, the wave function\ differs only by a
constant.
 We shall say more about this 
in the next section.

We may therefore define a $\gamma$-based theory in the following way. 
The position
of the particle at time $t$, $Q_t$, is a point on the circle. The wave function\
is a smooth function on $\ensuremath{\mathbb{R}^{}}$ satisfying \eqref{ie:sper}. The
evolution of the wave function\ is governed by Schr\"odinger's equation\ on \ensuremath{\mathbb{R}^{}}. The
generated Bohmian velocity field projects, in a single-valued way, to
the circle. The position then evolves according to that velocity field. 

As for the relation to the Aharonov-Bohm setup, we make the following
claim, which will be established in more generality in
Section {\ref{iss:bunequ}}. Let $\beta \ensuremath{:=} \int_{D^2} \boldsymbol{B}
\boldsymbol{\omega}$. This integral requires a choice of orientation
and the choice induces an orientation on the circle. If increasing
$\theta$ on $\ensuremath{\mathbb{R}^{}}$ agrees with this orientation, then $\gamma
\ensuremath{:=} e^{-i \beta}$ will give the appropriate
theory.


The above story holds much more generally. We first state the
following principle for scalar-valued wave functions. We then explain the
terminology and notation in a following subsection. An even more
general statement is given in the next section.
\begin{principle}[Abelian Quantization Principle, covering formulation, scalar]
To specify a Bohmian theory involving complex-valued wave functions, 
one must choose  a Riemannian configuration space
 \ensuremath{\mathcal{Q}}, a potential $V$, and a 
 character \ensuremath{\gamma}\
 of the fundamental group \ensuremath{\pi_1 ({\ensuremath{\mathcal{Q}}})}\ of \ensuremath{\mathcal{Q}}.
The Bohmian
theory corresponding to the character \ensuremath{\gamma}\ is defined by
 \eqref{ie:rbe}, where $\ensuremath{\psi}_t$ is a smooth wave function\ on \ensuremath{\widehat{\ensuremath{\mathcal{Q}}}}\
 satisfying the character's associated periodicity condition 
 and locally satisfying \eqref{ie:rse},
 i.e.\ $\ensuremath{\psi}_t$ evolves according to the Schr\"odinger\ dynamics generated by the lift
 of $H$ from \ensuremath{\mathcal{Q}}\ to \ensuremath{\widehat{\ensuremath{\mathcal{Q}}}}. 
\end{principle}

\subsection{Notations and relevant facts}

The notations are summarized by
the table:
\begin{center}
\begin{tabular}{ccc}
Configuration space & \ensuremath{\mathcal{Q}} & $S^1$ \\
Universal covering space & \ensuremath{\widehat{\ensuremath{\mathcal{Q}}}} & \ensuremath{\mathbb{R}^{}} \\
Points in \ensuremath{\mathcal{Q}}, \ensuremath{\widehat{\ensuremath{\mathcal{Q}}}} & \ensuremath{q}, \ensuremath{\hat{q}} &  $e^{i\theta}$, $\theta$ \\
Projection map & $\ensuremath{\ensuremath{\rho}} : \ensuremath{\widehat{\ensuremath{\mathcal{Q}}}} \to \ensuremath{\mathcal{Q}}$ & $e^{i \cdot}: \ensuremath{\mathbb{R}^{}}
\to S^1$\\
Covering fiber over \ensuremath{q}, $e^{i\theta}$ & $\ensuremath{\ensuremath{\rho}}^{-1} (\ensuremath{q})$ &
$\{\theta + 2k\pi\}_{k \in \mathbb{Z}}$ \\
Fundamental group of \ensuremath{\mathcal{Q}} & \ensuremath{\pi_1 ({\ensuremath{\mathcal{Q}}})} & $\mathbb{Z}$ \\
Covering transformation & \ensuremath{\sigma} & $\sigma_k : \theta \to \theta+2\pi k$ \\
Covering group & \ensuremath{Cov(\ensuremath{\widehat{\ensuremath{\mathcal{Q}}}}, \ensuremath{\mathcal{Q}})} & $\sigma_k$ \\
Character of fundamental group & $\ensuremath{\gamma}_{\ensuremath{\sigma}}$ & $\ensuremath{\gamma}_{k} \ensuremath{:=}
\ensuremath{\gamma}^{k}_1$  \\
Bundle, lifted bundle & $E$, $\ensuremath{\widehat{E}}$ & $S^1\times \ensuremath{\mathbb{C}}$,
$\ensuremath{\mathbb{R}^{}}\times \ensuremath{\mathbb{C}}$ \\
\end{tabular}
\end{center}


Throughout this chapter, \ensuremath{\mathcal{Q}}\ will be 
 a Riemannian manifold with metric $g$. Its universal covering
space will be \ensuremath{\widehat{\ensuremath{\mathcal{Q}}}}. We shall always take \ensuremath{\widehat{\ensuremath{\mathcal{Q}}}}\ to be endowed with the
lifted metric $\hat{g}$; this assures that it is canonically locally
  isometric
to \ensuremath{\mathcal{Q}}. More to the point, the differential operators are locally
identical. 

The universal covering space \ensuremath{\widehat{\ensuremath{\mathcal{Q}}}}\  of \ensuremath{\mathcal{Q}}\ is, by
construction,   a
simply connected space. The
\emph{covering fiber} for $\ensuremath{q}\ \in \ensuremath{\mathcal{Q}}$ is the set of points in \ensuremath{\widehat{\ensuremath{\mathcal{Q}}}}\ that
project to \ensuremath{q}\ under the projection, or covering,  map, 
$\ensuremath{\ensuremath{\rho}}: \ensuremath{\widehat{\ensuremath{\mathcal{Q}}}} \to \ensuremath{\mathcal{Q}}$,  associated with 
the covering 
space. A function $f: \ensuremath{\widehat{\ensuremath{\mathcal{Q}}}} \to \ensuremath{\mathbb{C}}$ is projectable iff  $f(\ensuremath{\hat{q}}) =
f(\hat{r})$ whenever $\ensuremath{\ensuremath{\rho}}(\ensuremath{\hat{q}}) = \ensuremath{\ensuremath{\rho}}(\hat{r})$. 
The projection  $\tilde{f} :
\ensuremath{\mathcal{Q}} \to \ensuremath{\mathbb{C}}$ is defined by $\tilde{f}(\ensuremath{q}) = f(\ensuremath{\ensuremath{\rho}}^{-1}
(\ensuremath{q}))$.
Projectability of a vector field 
can be translated into demanding that  
$((D\ensuremath{\ensuremath{\rho}}) v^{\ensuremath{\psi}}) (\ensuremath{\hat{q}}) = ((D\ensuremath{\ensuremath{\rho}}) v^{\ensuremath{\psi}}) (\hat{r})$ whenever
$\ensuremath{\ensuremath{\rho}}(\ensuremath{\hat{q}}) = \ensuremath{\ensuremath{\rho}}(\hat r )$; $D\ensuremath{\ensuremath{\rho}}: T\ensuremath{\widehat{\ensuremath{\mathcal{Q}}}} \to T\ensuremath{\mathcal{Q}}$ is
the differential of the 
map \ensuremath{\ensuremath{\rho}}. The lift of a function or a vector field 
is always well-defined and it is the
inverse of projection. 

If $E$ is a bundle over \ensuremath{\mathcal{Q}}, then the lift of $E$, denoted
$\ensuremath{\widehat{E}}$, is a bundle over \ensuremath{\widehat{\ensuremath{\mathcal{Q}}}}; the fiber at \ensuremath{\hat{q}}\ is defined 
to be the fiber of $E$ at $\ensuremath{q}$, i.e.\ 
$\ensuremath{\widehat{E}}_{\ensuremath{\hat{q}}} \ensuremath{:=} E_{\ensuremath{q}}$ where $\ensuremath{q} \ensuremath{:=} \ensuremath{\ensuremath{\rho}}(\ensuremath{\hat{q}})$.
 When we do this, we can then lift sections of the vector bundle as
well as matrix-valued functions acting on the bundle. It is important
to realize that with this construction, it makes sense to ask if $v
\in E_{\ensuremath{\hat{q}}}$ is equal to $w \in E_{\hat{r}}$ whenever $\ensuremath{\hat{q}}$ and
$\hat{r}$ are elements of the same covering fiber.
As a particular example, the lift of the tangent
bundle of \ensuremath{\mathcal{Q}}\ to \ensuremath{\widehat{\ensuremath{\mathcal{Q}}}}\ is canonically isomorphic to 
the tangent bundle of \ensuremath{\widehat{\ensuremath{\mathcal{Q}}}}. 


A \emph{covering transformation} is an isometry which maps the covering
space to itself and which preserves the covering fiber s.  The group of such
transformations is the covering group and is denoted by \ensuremath{Cov(\ensuremath{\widehat{\ensuremath{\mathcal{Q}}}}, \ensuremath{\mathcal{Q}})}. 
The fundamental group at a point \ensuremath{q}, denoted
by \ensuremath{\pi_1 ({\ensuremath{\mathcal{Q}}, \ensuremath{q}})},
is the set of 
equivalence classes 
of closed loops, through \ensuremath{q}, where the equivalence relation is
that of homotopy, 
or smoothly deforming one curve into the other. 
The product in the group is concatenation; in particular, given
$\sigma \tau$, we first follow $\tau$ and then follow $\sigma$. This
is in contrast to the common definition of the product in the opposite
order. We do it this way as it seems more natural in terms of parallel
transport, which is of fundamental concern to us. 
The fundamental groups
at different points are isomorphic to each other as well as to the
covering group, but the isomorphisms are not canonical. 
By the fundamental group, written \ensuremath{\pi_1 ({\ensuremath{\mathcal{Q}}})}, we shall mean the
abstract group 
which the covering group and all of the fundamental groups are
isomorphic to.  In particular, a representation of the covering group
is a representation of the fundamental group. A character of the
fundamental group is any unitary, 1-dimensional representation of the
fundamental group, i.e., it is a homomorphism between the fundamental
group and the multiplicative group of the complex numbers of modulus
one. 

\section[AQP]{Bohmian mechanics\ and topology:\\ the Abelian Quantization Principle} \label{s:expo}


Since it is the same argument, we shall argue for the Abelian Quantization Principle\ in the
context of a 
 Bohmian theory in which the wave functions are sections of a vector bundle. We
start with a Riemannian configuration space $\ensuremath{\mathcal{Q}}$ and
 a vector bundle $E$ endowed with a connection and a parallel inner
product. We are also given a self-adjoint matrix-valued potential $V$
acting on $E$. We claim that, in addition to this information, we also
need to choose a character of the fundamental group to complete the
theory formation. That this is an \emph{a priori} necessity is
debatable since the theory corresponding to 
 the trivial character requires no need to
appreciate the other possibilities. But the physics certainly seems to
require an appreciation of these options.

As in the case of the circle, the idea is to appeal to the universal covering
space \ensuremath{\widehat{\ensuremath{\mathcal{Q}}}}\ in order to trivialize the topology. We lift all of the
objects of our theory to \ensuremath{\widehat{\ensuremath{\mathcal{Q}}}}. This includes lifting $E$ to
$\ensuremath{\widehat{E}}$
by defining $E_{\ensuremath{\hat{q}}} \ensuremath{:=} E_{\ensuremath{q}}$. By lifting the local
objects, we can write down a Schr\"odinger's equation\ for sections of $\ensuremath{\widehat{E}}$; we refer
to the Hamiltonian on \ensuremath{\widehat{\ensuremath{\mathcal{Q}}}}\ as the lifted Hamiltonian.  
 As before, we are only interested in wave functions which
generate a Bohmian evolution on \ensuremath{\mathcal{Q}}. In particular, the velocity
field defined in \eqref{ie:vbv}\ needs to be projectable. 

Although we no longer have an $S$, as we are dealing with vector
bundles, it is quite clear that projectability is guaranteed 
if we have a scalar periodicity condition as in the case of the
circle. More precisely, if $\ensuremath{\psi}$ is a section of $\ensuremath{\widehat{E}}$ satisfying 
\begin{equation} \label{e:percon}
\ensuremath{\psi} (\ensuremath{\sigma} \ensuremath{\hat{q}})
= \ensuremath{\gamma_{\ensuremath{\sigma}}} \ensuremath{\psi}(\ensuremath{\hat{q}})
\end{equation}
for every $\ensuremath{\hat{q}} \in \ensuremath{\widehat{\ensuremath{\mathcal{Q}}}}$ and $\ensuremath{\sigma} \in \ensuremath{Cov(\ensuremath{\widehat{\ensuremath{\mathcal{Q}}}}, \ensuremath{\mathcal{Q}})}$, and 
where \ensuremath{\gamma_{\ensuremath{\sigma}}}\ is a constant of modulus 
one, then it is clear that 
\[v^{\ensuremath{\psi}} (\ensuremath{\hat{q}}) = v^{\ensuremath{\psi}} (\ensuremath{\sigma}\ensuremath{\hat{q}}). 
\]

In order to say that we have a Bohmian dynamics on \ensuremath{\mathcal{Q}}, we
must demonstrate that if the initial wave function\ satisfies \eqref{e:percon},
then it remains projectable. In fact, we shall demonstrate that it
continues to satisfy the same periodicity condition. It follows from 
the linearity of Schr\"odinger's equation\  and the fact that the lifted Hamiltonian is locally the
same when compared at different points in the same covering fiber. 
The idea is that the wave function\ at different ``levels'' is the same
up to an irrelevant constant. Its evolution 
 should therefore be the same on the different levels and the
periodicity condition will thus be preserved. 

A more precise
argument is the following. If $\ensuremath{\psi}_t$ is any solution to Schr\"odinger's equation\ on
$\ensuremath{\widehat{E}}$,  then, for any covering transformation $\ensuremath{\sigma}$, we have
that $\ensuremath{\psi}_t^{\ensuremath{\sigma}} \ensuremath{:=} \ensuremath{\psi}_t \circ \ensuremath{\sigma}$ is another solution to
Schr\"odinger's equation.\footnote{Since the potential is identical at each point  of the
covering fiber, we only have to worry about the Laplacian. But since \ensuremath{\sigma}\
is an isometry we have $\ensuremath{\Delta} \ensuremath{\psi}_t^{\ensuremath{\sigma}} = (\ensuremath{\Delta} \ensuremath{\psi}_t) \circ
{\ensuremath{\sigma}}$.} On the other hand, if we multiply $\ensuremath{\psi}_t$ 
by a constant scalar, we still have a
solution. Using the periodicity condition and uniqueness of initial
conditions, we have that
the periodicity condition is preserved.\footnote{
Assume that $\ensuremath{\psi}_0$   satisfies the periodicity
condition. Let $\ensuremath{\psi}_t$ be the solution of Schr\"odinger's equation\ with that initial
condition. Define $\phi_t^{\ensuremath{\sigma}} \ensuremath{:=} \ensuremath{\gamma_{\ensuremath{\sigma}}} \ensuremath{\psi}_t$. By linearity, this is
a solution to Schr\"odinger's equation. As stated above, $\ensuremath{\psi}_t^{\ensuremath{\sigma}}$ is also a
solution. We claim that $\ensuremath{\psi}_0^{\ensuremath{\sigma}} = \phi_0^{\ensuremath{\sigma}}$. This is
clear since $\ensuremath{\psi}_0^{\ensuremath{\sigma}} (\ensuremath{\hat{q}})  \ensuremath{:=} \ensuremath{\psi}_0 (\ensuremath{\sigma}\ensuremath{\hat{q}}) = \ensuremath{\gamma_{\ensuremath{\sigma}}}
\ensuremath{\psi}_0(\ensuremath{\hat{q}}) \ensuremath{=:}  \phi_0^{\ensuremath{\sigma}} (\ensuremath{\hat{q}}) $. Thus, by uniqueness of
initial conditions, $\ensuremath{\psi}_t^{\ensuremath{\sigma}} = \phi_t^{\ensuremath{\sigma}}$ which implies
that the periodicity condition is preserved.}

Given the periodicity condition \eqref{e:percon}, it is a fact that
 the map $\gamma: \ensuremath{Cov(\ensuremath{\widehat{\ensuremath{\mathcal{Q}}}}, \ensuremath{\mathcal{Q}})} \to \ensuremath{\mathbb{C}}$,  must
be a  representation of
the covering group, as the following argument demonstrates. 
 Let $\ensuremath{\ensuremath{\sigma}_{1}}$, $\ensuremath{\ensuremath{\sigma}_{2}} \in  \ensuremath{Cov(\ensuremath{\widehat{\ensuremath{\mathcal{Q}}}}, \ensuremath{\mathcal{Q}})}$. Then
we have the following equalities
\begin{equation*} 
\ensuremath{\gamma_{\ensuremath{\ensuremath{\sigma}_{1}} \ensuremath{\ensuremath{\sigma}_{2}}}} \ensuremath{\psi} (\ensuremath{\hat{q}}) 
= \ensuremath{\psi}(\ensuremath{\ensuremath{\sigma}_{1}} \ensuremath{\ensuremath{\sigma}_{2}} \ensuremath{\hat{q}}) 
= \ensuremath{\gamma_{\ensuremath{\ensuremath{\sigma}_{1}}}} \ensuremath{\psi}(\ensuremath{\ensuremath{\sigma}_{2}} \ensuremath{\hat{q}})
= \ensuremath{\gamma_{\ensuremath{\ensuremath{\sigma}_{1}}}}\ensuremath{\gamma_{\ensuremath{\ensuremath{\sigma}_{2}}}} \ensuremath{\psi}(\ensuremath{\hat{q}}).
\end{equation*}
We therefore conclude the fundamental relation
 \begin{equation} \label{e:ccc}
\ensuremath{\gamma_{\ensuremath{\ensuremath{\sigma}_{1}} \ensuremath{\ensuremath{\sigma}_{2}}}} = \ensuremath{\gamma_{\ensuremath{\ensuremath{\sigma}_{1}}}}\ensuremath{\gamma_{\ensuremath{\ensuremath{\sigma}_{2}}}}.
\end{equation}
We thus have that  
each periodicity condition corresponds to a  character
 of the covering group. The periodicity
condition also forms a character of the fundamental group since the
fundamental group and covering group are isomorphic.


We have the following principle, which shall be denoted by \ensuremath{\text{AQP}_{\ensuremath{\widehat{\ensuremath{\mathcal{Q}}}}}}: 
\begin{principle}[Abelian Quantization Principle, \ensuremath{\widehat{\ensuremath{\mathcal{Q}}}}] 
To specify a Bohmian theory, one needs to pick a Riemannian manifold
 $\ensuremath{\mathcal{Q}}$, a 
 cc-Hermitian\ vector bundle $E$ over
 \ensuremath{\mathcal{Q}}, a potential $V$, and a character \ensuremath{\gamma}\
 of the fundamental group \ensuremath{\pi_1 ({\ensuremath{\mathcal{Q}}})}.
 The Bohmian
 theory $(\ensuremath{\mathcal{Q}}, E, V, \ensuremath{\gamma})_{\ensuremath{\widehat{\ensuremath{\mathcal{Q}}}}}$ is defined by 
\eqref{ie:ve}
where $v^{\ensuremath{\psi}_t}$ is the projection to \ensuremath{\mathcal{Q}}\ of the Bohmian velocity field
 defined on
 \ensuremath{\widehat{\ensuremath{\mathcal{Q}}}}\ and the wave function\
$\ensuremath{\psi}_t$ is a section of the lifted bundle $\ensuremath{\widehat{E}}$ satisfying
 \eqref{e:percon} and locally satisfying  \eqref{ie:vse}. 
\end{principle}

\subsection{Remarks}

\begin{enumerate}

\item The trivial periodicity condition, $\ensuremath{\sigma} \mapsto 1$ for all
\ensuremath{\sigma}, corresponds to lifted wave functions. Indeed, 
if we
lift a wave function\ from \ensuremath{\mathcal{Q}}\ to \ensuremath{\widehat{\ensuremath{\mathcal{Q}}}}, then, for every $\ensuremath{\hat{q}} \in
\ensuremath{\widehat{\ensuremath{\mathcal{Q}}}}$
and $\ensuremath{\sigma} \in \ensuremath{Cov(\ensuremath{\widehat{\ensuremath{\mathcal{Q}}}}, \ensuremath{\mathcal{Q}})}$,
the wave function\ satisfies the
periodicity condition
\begin{equation*}
\ensuremath{\psi} (\ensuremath{\sigma}\ensuremath{\hat{q}})
= \ensuremath{\psi}(\ensuremath{\hat{q}}).
\end{equation*}

\item Having a projectable Bohmian velocity field places no constraint
on the amplitude. 
One might then wonder whether the amplitude is allowed
to vary based on \ensuremath{\hat{q}}. Generically, such a condition will not be
preserved under the Schr\"odinger\ evolution. 
In other words, if $r_{\ensuremath{\sigma}} (\ensuremath{\hat{q}})$ is the proposed
amplitude, then in order to argue that it is preserved, we should have
$\ensuremath{\Delta} r_{\ensuremath{\sigma}} \ensuremath{\psi} = r_{\ensuremath{\sigma}} \ensuremath{\Delta}\ensuremath{\psi}$. That is to say, we need
$(\ensuremath{\Delta} r_{\ensuremath{\sigma}}) \ensuremath{\psi} + \nabla r_{\ensuremath{\sigma}} \cdot \nabla \ensuremath{\psi} =
0$. If $r_{\ensuremath{\sigma}}$ is independent of $\ensuremath{\psi}$, we see that
$\nabla r_{\ensuremath{\sigma}}$ must 
vanish and hence $r_{\ensuremath{\sigma}}$ is a constant.  

\item What if $r_{\ensuremath{\sigma}}$ is a constant? 
In standard quantum mechanics, the condition on wave functions on the covering
space is that their densities should project down unambiguously \cite{Mor92}. 
As this is not our motivation, one might wonder why \ensuremath{\gamma_{\ensuremath{\sigma}}}\ should be
of modulus one.  If the 
fundamental group is 
finite, or even just idempotent, then the constant must be a root of
unity as a result of $\gamma$ being a representation. If   
 the  group is not idempotent, then the representation properties do
not demand that \ensuremath{\gamma_{\ensuremath{\sigma}}}\ be 
of modulus one. But 
if \ensuremath{\gamma_{\ensuremath{\sigma}}}\ has a modulus not
equal to one, then $|\ensuremath{\psi}({\ensuremath{\sigma}^m} \ensuremath{\hat{q}})| = r_{\ensuremath{\sigma}}^{m}\ensuremath{\psi}(\ensuremath{\hat{q}})$,
for any integer $m$;
our amplitude is growing exponentially in space as we proceed from
level to level. 
With such a growth, there 
is reason to believe  that the evolution of the system
will be ill-defined. 

\item For vector bundles, multiplying by a scalar is not the only
possibility. We could imagine replacing $\gamma$ with a matrix,
$\Gamma$. Again, the problem is preservation of the condition under
the Schr\"odinger\ evolution. In particular, it is necessary that $[\Gamma, V]
= 0$. This could happen, e.g.\ $V$ is a scalar function times the
identity, but it should not be expected. 
 In section {\ref{s:bprin}}, we discuss
this possibility in the context of spin and demonstrate that it is not
possible. 

\item Another question is whether the periodicity condition
characterizes all projectable wave functions. This is clearly false since any
real wave function\ has a zero Bohmian velocity field and hence it projects
unambiguously. If we make the more important restriction that the wave function\
should be projectable for all time, it still fails. Indeed, it is
possible to have on the covering space a wave function\ which is essentially a real
eigenstate of the potential.
 It is therefore invariant and its
velocity field vanishes implying projectability. This is not a problem
as our goal is to formulate a theory. 
We are not interested in
having just one wave function, but we are rather demanding
 a robust class of wave functions that our
theory will have as initial conditions.\footnote{It can be argued  
 that a wave function\ which is projectable for all time, 
but does not satisfy the
periodicity condition, must have the property that the nodal set is
of codimension 1. Since wave functions generically have a nodal set of
codimension 2, as both the real and imaginary parts must vanish, we
can expect that the condition is only preserved under very special
conditions such as the wave function\ being an eigenstate and, therefore, not
evolving. This question was pursued in the context of identical
particles in \cite{Bac03}. As a particular example in that setup, and
one which should perhaps be read after the Chapter \ref{c:manpart},
consider 
 the hydrogen atom with two electrons and no interaction
between the electrons. Tensor the ground state with a real-valued excited
state. The state is then essentially invariant under the
evolution and has no particular symmetry. Since it is essentially
real, the Bohmian velocities are 
zero and, hence, consistent. We would expect, however, that under a
small perturbation of the potential (or including interactions) the
states would evolve in such a way as to lose their consistency
condition. 
} 

\item Different characters may give rise to equivalent theories. In the
examples we discuss, the theories are not equivalent, and this is
generically the case. 
If our bundles are 1-dimensional, then they are distinct. 
 We shall say more about this in Section {\ref{iss:beypri}} in a context for
which this can be more readily understood.

\end{enumerate}


\section{Different connections} \label{s:fb}

We have described a class of Bohmian theories that arise by considering
those wave functions,  defined on the covering space of the configuration
space, which provide an unambiguous evolution of the
configuration.
We motivated that line of pursuit with the circle in the
situation of the Aharonov-Bohm effect. In this section, we shall
pursue theories that arise by considering different connections, i.e.\
different possibilities for what the derivative is. This approach has
the advantage that it does not depend on Bohmian
considerations. 
Instead of changing the base space, we change the bundle
structures of the value space of the wave function. 
We automatically have a well-defined Bohmian velocity. 

\subsection{The circle}
We start with the circle. The usual wave functions over the circle are
sections of the trivial bundle $S^1 \times \ensuremath{\mathbb{C}}$. For trivial
bundles, there is a natural choice for a derivative; namely, form the
difference quotient. Since it is a trivial bundle, there is no issue
in comparing values in different fibers. But it is a fact that this
trivial connection is not the only flat connection on the circle. 
Indeed, any real-valued 1-form defines a new connection and if the
integral of the 1-form around a loop is not zero, then the new
connection is fundamentally different from the trivial connection. 

We claim that the set of complex numbers of modulus one are in
bijective correspondence with the distinct classes of connections.
Define
the real-valued 1-form  $d\theta$ on the circle such that its integral
around the circle is $2\pi$.\footnote{Parametrize the circle with the
multi-valued angle function $\theta$ and
differentiate. 
} Let
$\nabla$ be the trivial connection associated with
$S^1 \times \ensuremath{\mathbb{C}}$. Then $\nabla + i\frac{\beta}{2\pi}
d\theta$ defines a flat connection; the parallel transport of $v$
along the path 
$\alpha$ is given by
\begin{equation*}
\ensuremath{P_{{\alpha}}{v}} = e^{\int_{\alpha}i\frac{\beta}{2\pi}d\theta} v. 
\end{equation*}
Thus, around a closed loop of multiplicity $n$,
we have that multiplication by $e^{i n \beta}$ 
is the parallel transport operator. 
Note that although it is no longer the trivial flat bundle, the
bundle is still trivial as a vector bundle. 

As for the relation to the Aharonov-Bohm setup, the fact of the matter
is that $\beta$ 
should be the integral of the magnetic field over the
disk. Furthermore, there is a vector potential which generates the
magnetic field and, restricted to the circle, it is 
$\frac{\beta}{2\pi}d\theta$. 
This is essentially 
the original explanation of
the Aharonov-Bohm effect.

\subsection{General discussion}
The circle is a nice example since it implements the idea of changing
the connection on the bundle. In more general situations, this
changing of the connection actually requires a change of the
underlying bundle structure. We therefore formulate the general
discussion in terms of constructing a new bundle out of an old
bundle. This method provides the appropriate connection and it may or
may not change the underlying vector bundle. 


Inspired by the case of fermions, as discussed in the next chapter, 
we shall formulate the new principle slightly differently than one
might expect: tensor
the original bundle with a flat cc-Hermitian\ line bundle.
It is a fact that the flat cc-Hermitian\ line
bundles are in bijective correspondence, up to equivalences,
 with the characters of the
fundamental group. Indeed, 
the  representation associated with
a flat bundle is its holonomy representation; the \emph{holonomy
representation}, at \ensuremath{q}, is the map taking 
each element $\ensuremath{\sigma} \in \ensuremath{\pi_1 ({\ensuremath{\mathcal{Q}}, \ensuremath{q}})}$ to  the matrix
associated with  
parallel transporting vectors around a representative curve of \ensuremath{\sigma}. Flatness
implies this map is well-defined, i.e.\ parallel transport on flat
bundles only depends
upon the homotopy class of the curve.  For a line bundle, a matrix
is just a complex number; to preserve the inner product, it must be an
element of the unit circle. Therefore, every flat line bundle has an
associated character. 

The connection for the tensor product $F \otimes
E$ of two bundles is defined using the Leibniz rule:
\begin{equation*} \nabla^{F\otimes E}( f \otimes e) = (\nabla^F f)
\otimes e + f 
\otimes (\nabla^E  e).
\end{equation*} 
It can also be realized by tensoring the holonomy operators together.
Computing the  Laplacian,  we find  
\begin{equation*} 
\ensuremath{\Delta}^{F\otimes E} (f \otimes e) = (\ensuremath{\Delta}^F f) \otimes e + f \otimes (\ensuremath{\Delta}^E
e) + 2 (\nabla^F f) \cdot \otimes (\nabla^E e).
\end{equation*} 
Furthermore, 
it is an easy
fact that  in
computing the 
curvature, the cross terms cancel rather than add.\footnote{The
curvature, $\omega$,  is the antisymmetric part of the second derivative. In
particular, for a tensor product, we have
 \[\begin{split}  
D^2 (f\otimes e) (x, y) - D^2(f \otimes e)(y,x) = & D^2f(x,y) \otimes
e + Df(x) \otimes De(y) + f \otimes D^2e(x,y) + \\ &Df(y) \otimes
De(x)  
- (D^2f(y,x) \otimes
e + 
Dfy \otimes De(x) + \\
&f \otimes D^2e(y,x) + Df(x) \otimes
De(y) ). \end{split} \]
Thus, we have $\omega^{F \otimes E} = \omega^{F} \otimes \mathrm{Id}
+ \mathrm{Id} \otimes \omega^{E} $. In particular, if $F$ is a flat
line bundle, we can appropriately say that the curvature is unchanged
under tensoring with $F$. }
Thus, if $E$ has
curvature and $F$ is a flat line bundle,
 then $F \otimes E$ will have the same
curvature as $E$. Any potential $V$, defined on $E$, goes to
$\mathrm{Id} \otimes V$, which may reasonably be viewed as $V$ since
$F$ is one dimensional, i.e $\mathrm{Id} = 1$. 


Thus, we have the Abelian Quantization Principle\ based on changing the connection; we shall
denote it by \ensuremath{\text{AQP}_{\nabla}}.
\begin{principle}[Abelian Quantization Principle, $\nabla$]
To specify a Bohmian theory, one needs to pick a Riemannian manifold
 $\ensuremath{\mathcal{Q}}$, a 
 cc-Hermitian\ vector bundle $E$ over
 \ensuremath{\mathcal{Q}}, a potential $V$, and a character \ensuremath{\gamma}\
 of the fundamental group \ensuremath{\pi_1 ({\ensuremath{\mathcal{Q}}})}.
Let $F$ be the flat
 cc-Hermitian\ line bundle corresponding to \ensuremath{\gamma}. The Bohmian
 theory $(\ensuremath{\mathcal{Q}}, E, V, \ensuremath{\gamma})_{\nabla}$ is defined by 
\eqref{ie:ve}
where  the wave function\
$\ensuremath{\psi}_t$ is a section of the bundle $F\otimes E$. 
\end{principle}


\subsection{Equivalence of the two approaches} \label{iss:bunequ}

The bundle approach and the covering space approach are completely
equivalent as we shall now explain. We start with the bundle $E$ 
and denote its holonomy operator for the path $\alpha$ by
$P^E_{\alpha}$.\footnote{We should emphasize that we
are not assuming that the connection on $E$ is flat. Thus, the holonomy
operator depends on the path itself and not just the homotopy
class of the path.} 
We
now have two approaches for forming theories. One is to consider
various periodicity conditions, involving scalars, on the covering
space; the other is to tensor $E$ with the various flat line
bundles. Both are characterized by a character of the fundamental
group. We contend that either way produces equivalent theories for
equivalent characters.

Before we produce the equivalence, we first discuss the 
canonical isomorphism
between the covering group and the fundamental group based at a point. There
are two natural bijections and it is the less obvious one which is
relevant, as we now explain. 
Let \ensuremath{\ensuremath{\rho}}\ be the projection  map
from \ensuremath{\widehat{\ensuremath{\mathcal{Q}}}}\ to \ensuremath{\mathcal{Q}}. Then $\ensuremath{\ensuremath{\rho}}(\alpha)  = \ensuremath{\ensuremath{\rho}}(\ensuremath{\sigma}
\alpha)$ for all paths $\alpha$ on \ensuremath{\widehat{\ensuremath{\mathcal{Q}}}}\ and for all $\ensuremath{\sigma} \in
\ensuremath{Cov(\ensuremath{\widehat{\ensuremath{\mathcal{Q}}}}, \ensuremath{\mathcal{Q}})}$. This is trivially true since the covering transformations map
each covering fiber\ to itself and \ensuremath{\ensuremath{\rho}}\ projects all points of the fiber to
the same base point. Given $\ensuremath{\hat{q}} \in
\ensuremath{\widehat{\ensuremath{\mathcal{Q}}}}$, and $\ensuremath{q} \ensuremath{:=} \ensuremath{\ensuremath{\rho}}(\ensuremath{\hat{q}})$, we can define
an isomorphism between \ensuremath{\pi_1 ({\ensuremath{\mathcal{Q}}, \ensuremath{q}})}\ and \ensuremath{Cov(\ensuremath{\widehat{\ensuremath{\mathcal{Q}}}}, \ensuremath{\mathcal{Q}})}. The obvious
bijection from $\ensuremath{Cov(\ensuremath{\widehat{\ensuremath{\mathcal{Q}}}}, \ensuremath{\mathcal{Q}})} \to \ensuremath{\pi_1 ({\ensuremath{\mathcal{Q}}, \ensuremath{q}})}$ 
is to map the covering transformation \ensuremath{\sigma}\ into the
projection of any path taking $(\ensuremath{\hat{q}} \to \ensuremath{\sigma} \ensuremath{\hat{q}})$. This is fine
as a bijection, but it fails to be an isomorphism.\footnote{To be an
  isomorphism, we would need $\ensuremath{\ensuremath{\rho}}(\ensuremath{\hat{q}} \to \ensuremath{\sigma} \tau \ensuremath{\hat{q}}) \ensuremath{\sim}
  \ensuremath{\ensuremath{\rho}}(\ensuremath{\hat{q}} \to \ensuremath{\sigma} \ensuremath{\hat{q}}) \circ \ensuremath{\ensuremath{\rho}}(\ensuremath{\hat{q}} \to \tau
  \ensuremath{\hat{q}})$. We certainly have $\ensuremath{\ensuremath{\rho}}(\ensuremath{\hat{q}} \to \ensuremath{\sigma} \tau \ensuremath{\hat{q}}) \ensuremath{\sim}
  \ensuremath{\ensuremath{\rho}}(\tau \ensuremath{\hat{q}} \to \ensuremath{\sigma} \tau \ensuremath{\hat{q}}) \circ \ensuremath{\ensuremath{\rho}}(\ensuremath{\hat{q}} \to \tau
  \ensuremath{\hat{q}})$ and therefore only need
   $\ensuremath{\ensuremath{\rho}}(\ensuremath{\hat{q}} \to \ensuremath{\sigma} \ensuremath{\hat{q}})
  \ensuremath{\sim} \ensuremath{\ensuremath{\rho}}(\tau \ensuremath{\hat{q}} \to \ensuremath{\sigma} \tau \ensuremath{\hat{q}})$. But, in actual fact,
  we can only conclude $\ensuremath{\ensuremath{\rho}}(\ensuremath{\hat{q}} \to \ensuremath{\sigma} \ensuremath{\hat{q}})
  \ensuremath{\sim} \ensuremath{\ensuremath{\rho}}(\tau \ensuremath{\hat{q}} \to \tau \ensuremath{\sigma} \ensuremath{\hat{q}})$ by applying $\tau$ to
  the left hand side. If the covering group is not abelian, then this
bijection is not an isomorphism, but an anti-isomorphism.} 
To find an isomorphism, we map it to the
inverse of the above map. More precisely, map \ensuremath{\sigma}\ into
$\ensuremath{\ensuremath{\rho}}(\ensuremath{\sigma} \ensuremath{\hat{q}} \to \ensuremath{\hat{q}}) \ensuremath{\sim} \ensuremath{\ensuremath{\rho}}(\ensuremath{\hat{q}} \to \ensuremath{\sigma}^{-1}
\ensuremath{\hat{q}})$. This is an isomorphism as is easy to check.\footnote{To be an
  isomorphism, we would need $\ensuremath{\ensuremath{\rho}}(\ensuremath{\hat{q}} \to (\ensuremath{\sigma} \tau)^{-1} \ensuremath{\hat{q}}) \ensuremath{\sim}
  \ensuremath{\ensuremath{\rho}}(\ensuremath{\hat{q}} \to \ensuremath{\sigma}^{-1} \ensuremath{\hat{q}}) \circ \ensuremath{\ensuremath{\rho}}(\ensuremath{\hat{q}} \to \tau^{-1}
  \ensuremath{\hat{q}})$. Since $\ensuremath{\ensuremath{\rho}}(\ensuremath{\hat{q}} \to (\ensuremath{\sigma} \tau)^{-1} \ensuremath{\hat{q}}) \ensuremath{\sim}
  \ensuremath{\ensuremath{\rho}}(\tau^{-1} \ensuremath{\hat{q}} \to \tau^{-1} \ensuremath{\sigma}^{-1} \ensuremath{\hat{q}}) \circ
\ensuremath{\ensuremath{\rho}}(\ensuremath{\hat{q}} \to \tau^{-1}
  \ensuremath{\hat{q}}),$
 we need 
$\ensuremath{\ensuremath{\rho}}(\ensuremath{\hat{q}} \to \ensuremath{\sigma}^{-1} \ensuremath{\hat{q}})
  \ensuremath{\sim} \ensuremath{\ensuremath{\rho}}(\tau^{-1} \ensuremath{\hat{q}} \to \tau^{-1} \ensuremath{\sigma}^{-1} \ensuremath{\hat{q}})$ which
  is what we have 
   by applying $\tau^{-1}$ to
  the left hand side.} Hence, we have a representation of the
fundamental group if and only if the image under this isomorphism is a
representation of the covering group. Abstractly, they will be the
same representation.  This isomorphism is exactly what comes out of
the following equivalence. 


\begin{theorem} 
The principles generate equivalent Bohmian theories, i.e.\ 
\[(\ensuremath{\mathcal{Q}}, E, V, \ensuremath{\gamma})_{\ensuremath{\widehat{\ensuremath{\mathcal{Q}}}}} \ensuremath{\simeq} (\ensuremath{\mathcal{Q}}, E, V,
\ensuremath{\gamma})_{\nabla}.\] 
\end{theorem}


\begin{proof}
Let $F$ be the flat line bundle associated with \ensuremath{\gamma}. 
The proof will go as follows. We start with $(\ensuremath{\mathcal{Q}}, E, V,
\ensuremath{\gamma})_{\nabla}$. We 
lift the bundle $G \ensuremath{:=} F\otimes E$ and show that the lifted sections are
essentially the sections of the lift of $E$ which satisfy the
periodicity condition \ensuremath{\gamma}. We also show that this mapping commutes
with 
the evolutions. 


We start by lifting the bundle $F$ to
the covering 
space; as usual, we denote its lift by $\ensuremath{\widehat{F}}$. 
Since $\ensuremath{\widehat{F}}$ is flat and \ensuremath{\widehat{\ensuremath{\mathcal{Q}}}}\ is simply connected, 
it is trivializable. Let $\phi$ be a trivialization map of
$\ensuremath{\widehat{F}}$. That is to say, $\phi: \ensuremath{\widehat{F}} \to \ensuremath{\widehat{\ensuremath{\mathcal{Q}}}} \times \ensuremath{\mathbb{C}}$
should be a bundle isomorphism preserving the connection, i.e.\
$\nabla(\phi( \ensuremath{\psi})) = \phi( \nabla\ensuremath{\psi})$ for any section \ensuremath{\psi}\ of $\ensuremath{\widehat{F}}$.  
We now lift $G$ to $\ensuremath{\widehat{G}} =\ensuremath{\widehat{F}} \otimes\ensuremath{\widehat{E}}$. 
The Bohmian theory on $\ensuremath{\widehat{G}}$ with the trivial periodicity
condition is completely equivalent to the Bohmian theory on $G$. 
Apply $\Phi \ensuremath{:=} \phi \otimes \text{Id}$ to $\ensuremath{\widehat{G}}$; this is
clearly an isomorphism between $\ensuremath{\widehat{G}}$ and $\ensuremath{\widehat{E}}$. Furthermore, it
commutes with the covariant derivative due to $\phi$ being
connection-preserving. We also know that it commutes with 
 all potentials since it is essentially a scalar.  
More to the point, if $\ensuremath{\psi}_t$ is a section of
$\ensuremath{\widehat{G}}$ and it is a solution to the appropriate
Schr\"odinger's equation, then we have
\begin{align}
i\ensuremath{\frac{\partial {\ensuremath{\psi}_t}}{\partial {t}}} &= H \ensuremath{\psi}_t  \Leftrightarrow i\ensuremath{\frac{\partial {\Phi(\ensuremath{\psi})_t}}{\partial {t}}} = H
\Phi(\ensuremath{\psi})_t  \\
v^{\ensuremath{\psi}_t} &= v^{\Phi(\ensuremath{\psi})_t}.
\end{align}
Thus, once we establish the claim below, we will have shown that these
theories are equivalent.

Our claim is
that $\ensuremath{\psi}_G$ is a lifted section of $G$ if and only if $\ensuremath{\psi}_E \ensuremath{:=}
\Phi(\ensuremath{\psi}_G)$ satisfies $\ensuremath{\psi}_E (\ensuremath{\sigma} \ensuremath{\hat{q}}) = \ensuremath{\gamma}_{\ensuremath{\sigma}}\ensuremath{\psi}_E
(\ensuremath{\hat{q}}) $. 
To see this, consider what the partial trivialization
  is. Given points $\ensuremath{\hat{q}}, \ensuremath{\widehat{r}} \in \ensuremath{\widehat{\ensuremath{\mathcal{Q}}}}$, then the
  trivialization $\phi$ satisfies 
\[\phi_{\ensuremath{\widehat{r}}} = P_{\ensuremath{\hat{q}} \to \ensuremath{\widehat{r}} }^{\ensuremath{\widehat{\ensuremath{\mathcal{Q}}}} \times \ensuremath{\mathbb{C}}} \phi_{\ensuremath{\hat{q}}}
  P_{\ensuremath{\widehat{r}} \to \ensuremath{\hat{q}}}^{\ensuremath{\widehat{F}}} = \phi_{\ensuremath{\hat{q}}}
  P_{\ensuremath{\widehat{r}} \to \ensuremath{\hat{q}}}^{\ensuremath{\widehat{F}}};\] 
 this follows from the fact that the trivialization is parallel
    and that parallel transport on the trivial flat bundle is
    trivial. Thus, if we  map  $\ensuremath{\psi}_G$, a lifted section of $G$, using this
    partial trivialization to a section of $\ensuremath{\widehat{E}}$, $\ensuremath{\psi}_E$, we find the
    following:
\begin{equation*}
\begin{split}
 \ensuremath{\psi}_E (\ensuremath{\sigma} \ensuremath{\hat{q}}) = & \Phi_{\ensuremath{\sigma} \ensuremath{\hat{q}}} \ensuremath{\psi}_G (\ensuremath{\sigma} \ensuremath{\hat{q}}) =
\Phi_{\ensuremath{\sigma} \ensuremath{\hat{q}}} \ensuremath{\psi}_G ( \ensuremath{\hat{q}})  \\
= & \Phi_{\ensuremath{\hat{q}}} (P_{\ensuremath{\sigma} \ensuremath{\hat{q}} \to
  \ensuremath{\hat{q}}}^{\ensuremath{\widehat{F}}}\otimes \text{Id})
 \ensuremath{\psi}_G(\ensuremath{\hat{q}}) = \Phi_{\ensuremath{\hat{q}}} (P_{\ensuremath{\sigma} \ensuremath{\hat{q}} \to
  \ensuremath{\hat{q}}}^{\ensuremath{\widehat{F}}} \otimes \text{Id}) \Phi_{\ensuremath{\hat{q}}}^{-1} 
\ensuremath{\psi}_E(\ensuremath{\hat{q}}) = \gamma_{\ensuremath{\ensuremath{\rho}}(\ensuremath{\sigma} \ensuremath{\hat{q}} \to  \ensuremath{\hat{q}})} \ensuremath{\psi}_E(\ensuremath{\hat{q}});
\end{split}
\end{equation*}
it is important to note that $P_{\ensuremath{\sigma} \ensuremath{\hat{q}} \to
  \ensuremath{\hat{q}}}^{\ensuremath{\widehat{F}}} = \gamma_{\ensuremath{\ensuremath{\rho}}(\ensuremath{\sigma} \ensuremath{\hat{q}} \to  \ensuremath{\hat{q}})}$ which is
  the defining property of $F$. This demonstrates that the mapping
  injects into the set of wave functions in $\ensuremath{\widehat{E}}$ 
satisfying the periodicity condition. We must show that it is
  onto. But this is just a rewriting of the above series of
  equations.


\end{proof}

As an immediate application of this theorem, we have that the
covering space description of the Aharonov-Bohm effect is equivalent
to the vector potential description; it is important to appreciate the
correct 
isomorphism in order to make the correct correspondence. 


\section{Self-adjoint extensions from cutting} \label{s:sa}


We now pursue the functional analytic approach. This approach is
concerned with different self-adjoint extensions of the formal
Hamiltonian on the same \ensuremath{L^{2}}\ space. We will formulate a principle
based on this approach and it will correspond to the previous approaches. In
particular, for each character of the fundamental group, we shall have
a corresponding Bohmian theory. 
We first introduce the rough picture and then provide bite-sized
pieces of a precise treatment.

The idea is to cut \ensuremath{\mathcal{Q}}\ in order to make it simply connected. That
is to say, we shall specify a set $\ensuremath{\kappa} \subset \ensuremath{\mathcal{Q}}$ such that $
\ensuremath{\bar{\ensuremath{\mathcal{Q}}}_{\ensuremath{\kappa}}} \ensuremath{:=} \ensuremath{\mathcal{Q}} \setminus \ensuremath{\kappa}$ is a simply connected open
set. We shall also demand that \ensuremath{\kappa}\ be appropriately smooth and
minimal. If $E$ is a bundle over \ensuremath{\mathcal{Q}}, then $\ensuremath{\bar{E}_{\ensuremath{\kappa}}}$ denotes the
restriction of  $E$ to $\ensuremath{\bar{\ensuremath{\mathcal{Q}}}_{\ensuremath{\kappa}}}$. 
We also have an induced Hamiltonian  acting on
the sections with compact support in \ensuremath{\bar{\ensuremath{\mathcal{Q}}}_{\ensuremath{\kappa}}}. But  such
wave functions no longer form a domain of essential self-adjointness for $H$ in the
Hilbert space $\ensuremath{L^{2}}(\ensuremath{\mathcal{Q}}, E)$. We
therefore need to choose a self-adjoint extension of $H$. 
Generally, there are many different extensions. 
For us, a self-adjoint extension will be acceptable only if 
the resulting Bohmian evolution cannot be used to detect which
cut was used to generate the theory. 


\subsection{The circle}

 
Let us consider our canonical example, $\ensuremath{\mathcal{Q}} = S^1$, $E=
\ensuremath{\mathcal{Q}} \times \ensuremath{\mathbb{C}}$. We parametrize the circle with 
the angle $\theta$.
A cut \ensuremath{\kappa}\ corresponds to choosing a point; say $\ensuremath{\kappa} = 0$ in the
parametrization.   
Once we cut it, then we may view the space as the interval $[0,
 2\pi]$. We wish to consider various self-adjoint extensions of the
Laplacian defined on $\ensuremath{C^{\infty}_0}([0,2\pi])$. The domains of essential
self-adjointness which correspond to the ones we have found before
consist of functions in 
$C^{\infty} ( [0, 2\pi])$
 satisfying the boundary condition $D^n \ensuremath{\psi}( 2 \pi) = \gamma
D^n \ensuremath{\psi} (0)$.\footnote{Here $\gamma$ must be of modulus one so that the
Hamiltonian is symmetric on this domain as an integration by
parts argument demonstrates.}  
This corresponds to the
periodicity condition $\ensuremath{\psi} (\theta + 2 \pi) = \gamma
 \ensuremath{\psi} (\theta)$ on \ensuremath{\widehat{\ensuremath{\mathcal{Q}}}}\ 
and it also corresponds to having 
the connection $\nabla -  \frac{\ln(\gamma)}{2 \pi} d\theta$ on
$ \ensuremath{\mathcal{Q}}\times \ensuremath{\mathbb{C}}$.


It is a  fact that the full set of
self-adjoint extensions  is a family parameterized by four
parameters. The extensions above are
 characterized by just one parameter. How do we discard the other
extensions? The idea is that the Bohmian evolution should not depend
upon which cut was chosen; in this regard, a cut should be thought of
as a choice of coordinates. For example, if we choose the extension
with boundary conditions that \ensuremath{\psi}\ vanishes at both endpoints, then
we know that typically, the particle's trajectory will not cross the
cut point. This is the kind of dependence that we are trying to
avoid.  
A variety of extensions have the property that the velocity
field is ill-defined at the cut point; we certainly would exclude
those extensions. There is one family of
extensions for which the velocity field is defined and continuous at
the cut point, but it is not smooth
at the cut point. We exclude them as well. This leaves us with only
the ones we had previously.

We have roughly described why the other cuts are unsuitable from our
perspective. The question is now whether the boundary conditions we
consider good actually satisfy being cut-independent. 
Choose $\gamma_1$ and $\gamma_2$ such that  $\gamma =
\frac{\gamma_1}{\gamma_2}$.  
Another cut, $\ensuremath{\kappa'} \in (0, 2\pi)$, would correspond to
$[\ensuremath{\kappa'}, \ensuremath{\kappa'}
+2\pi]$. A unitary equivalence map \ensuremath{U^{\ensuremath{\kappa}{'}  \ensuremath{\kappa}{}}} from \ensuremath{\kappa'}\ to \ensuremath{\kappa}, which
preserves the Bohmian velocities except at the cut points, is
\[
\ensuremath{U^{\ensuremath{\kappa}{'}  \ensuremath{\kappa}{}}}(\theta) \ensuremath{\psi}(\theta) \ensuremath{:=} \begin{cases}
                         \gamma_1 \ensuremath{\psi}(\theta) & \theta \in (0, \ensuremath{\kappa'})\\
                         \gamma_2 \ensuremath{\psi}(\theta) & \theta \in (\ensuremath{\kappa'}, 2\pi)
                        \end{cases}.
\]
As this map is constant away from the cuts, it is rather
self-evident that the Bohmian velocity field is unchanged under this mapping.
Furthermore, this is a bijection between the domains; choosing the
$\gamma_i$ so that  the ratio is $\gamma$ is the key fact. This
ensures that the 
discontinuity at \ensuremath{\kappa'}\ disappears while the smoothness at
\ensuremath{\kappa}\ is replaced with a discontinuity
 ensuring that the wave function\ satisfies the
boundary condition.
As the wave function\ evolution will commute with this mapping, we can see that
the Bohmian evolution is independent of which cut is chosen in the
formation of the Hamiltonian.


Let us make a couple of remarks about approaches that did not work. 
If we  were to demand that the wave function's  evolution
is independent of the cut, then the only allowed boundary condition is
 the trivial periodic
boundary condition;  the only domain which is
independent of the cut is the one in which \ensuremath{\psi}, and its derivatives,
agree at the boundary.
 Another possibility, one which is too weak, is
that of unitary equivalence. On the circle, given any extension based
on a cut, we can implement a rotation on the circle and obtain a
unitarily equivalent evolution based on another cut. That is to say,
unitary equivalence excludes nothing. 


\subsection{What is a cut?}
We first need to define what a cut is. After doing that, we 
consider certain illuminating examples although the most illuminating
is in the context of identical particles as discussed in the next chapter. 

A cut, \ensuremath{\kappa}, is a closed subset of \ensuremath{\mathcal{Q}}\ satisfying
\begin{enumerate}
\item Its complement, \ensuremath{\bar{\ensuremath{\mathcal{Q}}}_{\ensuremath{\kappa}}}, is a simply connected
set.\footnote{We take simply connected to also imply that the space is
connected. 
}
\item If $\ensuremath{\bar{\ensuremath{\mathcal{Q}}}_{\ensuremath{\kappa}}} \subset A$ and $A$ is an open, simply connected
subset of \ensuremath{\mathcal{Q}}, then $A = \ensuremath{\bar{\ensuremath{\mathcal{Q}}}_{\ensuremath{\kappa}}}$.
\item The cut \ensuremath{\kappa}\ should be smooth in the sense that about any point
$p \in \ensuremath{\kappa}$, there is a neighborhood $V$ such that $V \cap \ensuremath{\kappa}$ is
the union of a finite number of regular, $(n-1)$-dimensional
submanifolds.\footnote{A regular
submanifold is one in which every point of it has a coordinate
chart in the manifold 
in which the submanifold looks like a subspace in that coordinate
system. This assumption of regularity 
is to avoid having a boundary such as the following in
the plane:  
$\{(x, \sin(\frac{1}{x}))| 0<x\leq 1\} \cup \{(0,y)| y \in [-1,1]\}$.
 That is not the kind of boundary we want. Note that in that
example, any small neighborhood about the origin has an
infinite number of disconnected  lines arising from 
the intersection of the submanifold
with the neighborhood.}   
\end{enumerate}


For $S^1$, an $(n-1)$-dimensional submanifold is a point. If we remove
more than one point, then the resulting space ceases to be
connected. Thus, a cut \ensuremath{\kappa}\ for the circle is the removal of a single
point. For the torus, $T^2$, an example of a cut is the removal of a
horizontal circle and a vertical circle. Their intersection is
typically a
single point. In the usual representation of the torus as a square in
the plane, the boundary of the square is the cut \ensuremath{\kappa}\ and the
interior is the fundamental domain, \ensuremath{\bar{\ensuremath{\mathcal{Q}}}_{\ensuremath{\kappa}}}. Analogously,
if we take a cube in \ensuremath{\mathbb{R}^{3}}\ and identify opposite sides to obtain
$T^3 = S^1\times S^1 \times S^1$, then the
boundary of the cube is a cut. Notice that it consists of the union
of three copies of $T^2$. The twelve edges of the cube are the
three circles each of which  corresponds to an intersection of two of these
surfaces. The corners
 correspond to the single point of $T^3$ in which all three
manifolds intersect.  In contrast, a vertical cylinder's cut is a vertical
line which is a single $(2-1)$-dimensional submanifold with no
self-intersections.  


\subsection{BV-equivalence and cut-independent self-adjoint extensions}


We shall now make a precise expression as to what it means for the
 Bohmian evolution to be independent of the cut. 
We define two self-adjoint operators $H_1$  and $H_2$ to be
\emph{BV-equivalent} iff there exists a unitary operator $U: 
L^2[\ensuremath{\mathcal{Q}},E] \to L^2[\ensuremath{\mathcal{Q}}, E]$ such that 
\begin{enumerate}
\item  $v^{\ensuremath{\psi}} (\ensuremath{q}) 
= v^{U \ensuremath{\psi}} (\ensuremath{q})$ almost everywhere (w.r.t. Lebesgue measure)
\item $U H_1 \ensuremath{\psi} = H_2 U \ensuremath{\psi}$ for every $\ensuremath{\psi}$ in the domain of $H_1$. 
\end{enumerate}
This does define an equivalence relation on the set of self-adjoint
operators.\footnote{Composition of two BV-equivalences is clear
enough.  As for the inverse of BV-equivalence being a BV-equivalence,
the first property is still clear. For the
second property, the domain is the issue. To resolve it, take
 the adjoint of the equation and use the facts
that both $H_1$ and $H_2$ are self-adjoint and $U$ is unitary. Thus,
one can conclude $U^{-1} H_2 = H_1 U^{-1}$.} 
We shall state what these conditions imply, but we shall only give a
rough argument in establishing  the
implications. 


Since the velocity fields must be defined at the same places (almost
everywhere), the supports of $\ensuremath{\psi}$ and $U \ensuremath{\psi}$ must be
the same. This implies that $U$ is a unitary matrix-valued function over
\ensuremath{\mathcal{Q}}, i.e.\ $(U \ensuremath{\psi}) (\ensuremath{q})  = U(\ensuremath{q}) \ensuremath{\psi}(\ensuremath{q}) $, with $U(\ensuremath{q})$ a
unitary matrix acting on the value space of \ensuremath{\psi}.\footnote{This follows
from the fact that $U$ commutes with all of the position spectral
measures. If we were discussing only scalar valued wave functions, we could 
then say that $U$ is diagonalizable in the position
representation. The reason for the commutation is that if $P_A$ is a
projection operator onto the wave functions with support in $A$, then $U P_A
(\ensuremath{\psi}_A + \ensuremath{\psi}_{A^c}) = U \ensuremath{\psi}_A  = P_A( U\ensuremath{\psi}_A + U\ensuremath{\psi}_{A^c})$; all that
is necessary is that $U$ preserves supports.}
Furthermore, $U$
should be smooth except for a set of measure zero; the set of measure zero
arises since we only demand agreement up to a set of measure zero. We can
then use the 
preservation of the velocities to conclude that $U$ has to be 
constant 
in any region of smoothness.\footnote{We really mean parallel, i.e.\ $\nabla U =0$.  
We shall generally use the term constant especially since $U$ will
generally 
be a pointwise 
multiple of the identity.} 
The second condition implies that $U$ bijectively maps
the domain of $H_1$ to the domain of $H_2$.\footnote{From an orthodox
perspective, one would presumably demand that the local density is
unaffected by the unitary transformation. This immediately leads to
$U(q)$. To argue that it is locally constant, one would need to invoke
the commutation with the Hamiltonian. In particular, one needs 
$[\ensuremath{\Delta}, U] =0$. This immediately leads to, acting on a test function
\ensuremath{\psi}, $(\ensuremath{\Delta} U) \ensuremath{\psi} + \nabla U \cdot \nabla \ensuremath{\psi} = 0$. Thus, $U$ is
locally constant.}

Equipped with a satisfactory notion of equivalence, we shall state
precisely what  cut-independence of the extension means.
Given a cut \ensuremath{\kappa}, we define $\ensuremath{H_{\ensuremath{\kappa}}}^{0}$ to be the formal Hamiltonian 
 $ -\ensuremath{\Delta} +V$  with domain 
$\ensuremath{C^{\infty}_0}(\ensuremath{\bar{E}_{\ensuremath{\kappa}}})$. By $\ensuremath{H_{\ensuremath{\kappa}}}$, we shall mean a self-adjoint extension 
of $\ensuremath{H_{\ensuremath{\kappa}}}^{0}$. 
Then  \ensuremath{H_{\ensuremath{\kappa}}}\ is
\emph{cut-independent} iff for each cut \ensuremath{\kappa'}, there is a
corresponding $\ensuremath{H_{\ensuremath{\kappa'}}}$  which is BV-equivalent to \ensuremath{H_{\ensuremath{\kappa}}}. 
Although there are many choices, we shall simply assume that 
a 
BV-equivalence map taking $\ensuremath{H_{\ensuremath{\kappa}}} \to \ensuremath{H_{\ensuremath{\kappa'}}}$
 has
been chosen; it will be 
denoted by
\ensuremath{U^{\ensuremath{\kappa}{}  \ensuremath{\kappa}{'}}}. 


\subsection{The principle}

We now establish a family of cut-independent self-adjoint extensions,
one for each character. The construction is to appeal to the covering
space to obtain the relevant evolution.  
To do this, we  show that given a periodicity
condition on the covering space, we do have a corresponding
cut-independent self-adjoint extension.

In many ways, the cut approach is a non-smooth version of the
covering space approach. 
Indeed, select a fundamental domain\footnote{A fundamental domain is a
region of the covering space which bijectively corresponds, under the
projection map, to \ensuremath{\bar{\ensuremath{\mathcal{Q}}}_{\ensuremath{\kappa}}}. In more general contexts, the boundary of
a fundamental domain is allowed to be quite wild.} $R$ on the covering space
whose boundary is contained in the lift of \ensuremath{\kappa}. 
Given a character, \ensuremath{\gamma},
define
\[\Lambda_{\ensuremath{\gamma_{}}} \ensuremath{:=} \{\phi \in \ensuremath{\mathcal{D}}(H)
| \phi(\sigma
\ensuremath{\hat{q}}) = \ensuremath{\gamma_{\ensuremath{\sigma}}} \phi(\ensuremath{\hat{q}}) \forall \ensuremath{\hat{q}} \in \ensuremath{\widehat{\ensuremath{\mathcal{Q}}}}\},
\]
the set of smooth sections of the lifted bundle
$\ensuremath{\widehat{E}}$ satisfying the periodicity condition defined by \ensuremath{\gamma}. 
The unitary evolution on the covering space preserves this
domain and preserves the \ensuremath{L^{2}}\ norm of each section restricted to the
region. Define the dense domain of $\ensuremath{L^{2}}(\ensuremath{\mathcal{Q}}, E)$
\[
\ensuremath{E_{\ensuremath{\kappa}, \ensuremath{\gamma_{}}}}  \ensuremath{:=}
\{\ensuremath{\psi} | \exists \phi \in \Lambda_{\ensuremath{\gamma_{}}} \text{ s.t. }
 \forall \ensuremath{q} \in \ensuremath{\bar{\ensuremath{\mathcal{Q}}}_{\ensuremath{\kappa}}},  \ensuremath{\psi}(\ensuremath{q}) = \phi(\ensuremath{\hat{q}})  \text{ for
the unique } \ensuremath{\hat{q}} \in (\ensuremath{\ensuremath{\rho}}^{-1} (\ensuremath{q}) \cap R)\}
\]
and let $\ensuremath{\ensuremath{\rho}}_R(\phi)$ be the bijection making the correspondence between
the elements of $\Lambda_{\ensuremath{\gamma_{}}}$ and \ensuremath{E_{\ensuremath{\kappa}, \ensuremath{\gamma_{}}}}.
Then the covering space evolution defines a unitary evolution on
\ensuremath{E_{\ensuremath{\kappa}, \ensuremath{\gamma_{}}}}\ which takes the domain to itself. 
Thus, if we define \ensuremath{H_{\ensuremath{\kappa}}}\ to be the formal Hamiltonian defined on
\ensuremath{E_{\ensuremath{\kappa}, \ensuremath{\gamma_{}}}},  then \ensuremath{H_{\ensuremath{\kappa}}}\ is self-adjoint. 

To establish that \ensuremath{H_{\ensuremath{\kappa}}}\ is cut-independent, we choose another cut
\ensuremath{\kappa'}. We select a fundamental domain $R'$ on \ensuremath{\widehat{\ensuremath{\mathcal{Q}}}}\ 
 with a boundary formed
from the lift of \ensuremath{\kappa'}. We do the same process again to generate
\ensuremath{H_{\ensuremath{\kappa'}}}, choosing the same periodicity condition as the one for \ensuremath{H_{\ensuremath{\kappa}}}. 
The BV-equivalence mapping is extremely simple. Let $f^{RR'}$ be the
map taking $R$ to $R'$ which is defined away from the boundary.
Then the BV-equivalence map is
\[
\ensuremath{U^{\ensuremath{\kappa}{}  \ensuremath{\kappa}{'}}} \ensuremath{\psi}(\ensuremath{q}) 
= \ensuremath{U^{\ensuremath{\kappa}{}  \ensuremath{\kappa}{'}}} \ensuremath{\ensuremath{\rho}}_R(\phi) (\ensuremath{q}) 
\ensuremath{:=} \phi(f^{RR'}(\ensuremath{\ensuremath{\rho}}^{-1}_R(\ensuremath{q})) 
=\ensuremath{\gamma}_{\sigma(\ensuremath{q}, f^{RR'})} \ensuremath{\ensuremath{\rho}}_R(\phi) (\ensuremath{q}).
\]
In words, we take our wave function\ in \ensuremath{E_{\ensuremath{\kappa}, \ensuremath{\gamma_{}}}}, lift it to the appropriate wave function\ on
the covering space, using the level $R$, then we project down the
corresponding parts of that wave function\ defined over $R'$.  If
$\ensuremath{q} \in R \cap R'$, then the BV-equivalence map at \ensuremath{q}\ is the
identity. Elsewhere, it would be multiplication by the appropriate
phase induced by the periodicity condition.  
Thus, given a character and a cut, we have a 
 cut-independent self-adjoint extension of the formal 
Hamiltonian. 


As a general principle, we can state the following. 
\begin{principle}[Abelian Quantization Principle, H]
To specify a Bohmian theory, one needs to pick a Riemannian manifold
 $\ensuremath{\mathcal{Q}}$, a 
 cc-Hermitian\ vector bundle $E$ over
 \ensuremath{\mathcal{Q}}, a potential $V$, and a character \ensuremath{\gamma}\
 of the fundamental group \ensuremath{\pi_1 ({\ensuremath{\mathcal{Q}}})}.
 The Bohmian
 theory $(\ensuremath{\mathcal{Q}}, E, V, \ensuremath{\gamma})_H$ is defined by 
\eqref{ie:ve}
where  the wave function\
$\ensuremath{\psi}_t$ is a section of the bundle $ E$ and $H$ is the cut-independent
 self-adjoint extension defined by \ensuremath{\gamma}\ and the cut \ensuremath{\kappa}. 

\end{principle}


\subsection{ Characterization of cut-independence }
Given a cut-independent self-adjoint
extension whose BV-equivalence map is
of the form 
$ \gamma_{ \ensuremath{q} } \ensuremath{\text{Id}} $, 
then there is an associated character of \ensuremath{\pi_1 ({\ensuremath{\mathcal{Q}}})}.
In general, 
the various unitary representations of the fundamental group characterize the
possible cut-independent extensions; this can be argued in a similar
vein as below. As shall be explained in the next chapter, there is
good reason to ignore such possibilities. 

Therefore, for the rest of this section, BV-equivalence
maps will be multiples of the identity. 

\begin{theorem}
Let $E$ be a cc-Hermitian\ vector bundle over the Riemannian configuration
space \ensuremath{\mathcal{Q}}\ and $V$ a potential acting on $E$. 
Let \ensuremath{H_{\ensuremath{\kappa}}}\ be a cut-independent self-adjoint extension 
generating the Bohmian theory $(\ensuremath{\mathcal{Q}}, E, V)_{\ensuremath{H_{\ensuremath{\kappa}}}}$ whose
BV-equivalence maps are scalar multiples of the identity. Then there
exists a character $\ensuremath{\gamma}$ such that 
 $(\ensuremath{\mathcal{Q}}, E, V)_{\ensuremath{H_{\ensuremath{\kappa}}}} \ensuremath{\simeq} (\ensuremath{\mathcal{Q}}, E, V, \ensuremath{\gamma})_{\nabla}$.
\end{theorem}

From the equivalence proved in Section {\ref{s:fb}}, we can conclude that
the set of cut-independent self-adjoint extensions is characterized by
the cut and a character. 

\begin{proof}

We start with the cut-independent \ensuremath{H_{\ensuremath{\kappa}}}. 
Each pair $(\ensuremath{\bar{E}_{{\ensuremath{\kappa'}}}}, \ensuremath{U^{\ensuremath{\kappa}{}  \ensuremath{\kappa}{'}}})$ is to be a generalized chart for our
bundle.\footnote{By generalized, we mean that we are not actually
going to write down coordinate charts. Instead these are neighborhoods
with a bundle structure already implied. Our construction does not
require us to go to coordinates; it only requires us to have patches
that we can put together differently.} The transition
map  going from the chart of $\ensuremath{\kappa}'$ to the chart of $\ensuremath{\kappa}''$ is the
(BV-equivalence) map 
\[
\ensuremath{U^{\ensuremath{\kappa}{'}  \ensuremath{\kappa}{''}}} \ensuremath{:=} \ensuremath{U^{\ensuremath{\kappa}{}  \ensuremath{\kappa}{''}}}(\ensuremath{U^{\ensuremath{\kappa}{}  \ensuremath{\kappa}{'}}})^{-1} .
\] These maps are smooth away from
the cuts $\ensuremath{\kappa}'$ and $\ensuremath{\kappa}''$; we shall establish this later. 
We then form the disjoint union of all of these charts and define our vector
bundle to 
 be the set of equivalence classes under the relation
\begin{equation} \label{ei:eqvr}
(\ensuremath{q}, v)_{\ensuremath{\kappa}'} \ensuremath{\sim} (\ensuremath{q}, w)_{\ensuremath{\kappa}''} \text{ iff }
\ensuremath{U^{\ensuremath{\kappa}{'}  \ensuremath{\kappa}{''}}} v = w;
\end{equation}
By $(\ensuremath{q}, v)_{\ensuremath{\kappa}'}$ we mean that $\ensuremath{q} \in \ensuremath{\bar{\ensuremath{\mathcal{Q}}}_{{\ensuremath{\kappa'}}}}$
and $v \in E_{\ensuremath{q}}$. 
Because  BV-equivalence is an equivalence relation, the relation above
is indeed an equivalence relation.
Quite generally, 
it can be shown that this setup does define a bundle as is often done
in the context of physics. We shall call this bundle $G$. 
Parallel transport is defined, in a chart, as  the 
parallel transport in $E$.
Furthermore, \ensuremath{H_{\ensuremath{\kappa}}}\  defines a Schr\"odinger\ evolution
with the correct formal Hamiltonian on $G$; cut-independence is
important to show that the evolution in any of the other charts is
compatible. The first property for BV-equivalence establishes
that the Bohmian evolution will be the same either using \ensuremath{H_{\ensuremath{\kappa}}}\ or the
bundle $G$.  

All that remains is to demonstrate that $G = F_{\ensuremath{\gamma}} \otimes E$,
for some flat line bundle $F_{\ensuremath{\gamma}}$. For this, consider the
parallel transport of the element $v \in E_{\ensuremath{q}}$, $\ensuremath{q} \in
\ensuremath{\bar{\ensuremath{\mathcal{Q}}}_{\ensuremath{\kappa}}}$, around the closed loop
$\alpha$, parametrized in the interval $[0,1]$. 
For simplicity, assume that $\ensuremath{\bar{\ensuremath{\mathcal{Q}}}_{\ensuremath{\kappa}}}$ and
$\ensuremath{\bar{\ensuremath{\mathcal{Q}}}_{{\ensuremath{\kappa'}}}}$ form a cover of $\alpha$ and that it is possible to
find $t_1, t_2$
such that $\alpha_t \in \ensuremath{\bar{\ensuremath{\mathcal{Q}}}_{\ensuremath{\kappa}}}$ for $t \in [0,t_1] \cup [t_2, 1]$
while $\alpha_t \in \ensuremath{\bar{\ensuremath{\mathcal{Q}}}_{{\ensuremath{\kappa'}}}}$ for $t \in [t_1, t_2]$.  
Then parallel transport along $\alpha$ is 
\[
P_{\alpha_0 \to \alpha_1}^G v =
P_{\alpha_{t_2} \to \alpha_1}^{\ensuremath{\bar{E}_{{\ensuremath{\kappa}}}}} 
\ensuremath{U^{\ensuremath{\kappa}{'}  \ensuremath{\kappa}{}}}_{\alpha_{t_2}}  
P_{\alpha_{t_1} \to \alpha_{t_2}}^{\ensuremath{\bar{E}_{{\ensuremath{\kappa'}}}}}
\ensuremath{U^{\ensuremath{\kappa}{}  \ensuremath{\kappa}{'}}}_{\alpha_{t_1}} P_{\alpha_{0} \to \alpha_{t_1}}^{\ensuremath{\bar{E}_{{\ensuremath{\kappa}}}}} v
= \ensuremath{U^{\ensuremath{\kappa}{'}  \ensuremath{\kappa}{}}}_{\alpha_{t_2}} \ensuremath{U^{\ensuremath{\kappa}{}  \ensuremath{\kappa}{'}}}_{\alpha_{t_1}} P_{\alpha_0 \to
\alpha_1}^E v 
\]
where we used the fact that the BV-equivalence maps commute with the
parallel transport operators which is especially true as they are just
scalars. Thus, we have that parallel transport is just parallel
transport in $E$ times a phase factor. To finish the argument, we need
to argue that the phase factor depends only on the homotopy class of
the curve. This follows from two facts. First, within a chart, we can
smoothly vary the path $\alpha$ without changing the phase part since
BV-equivalence maps are constant away from the boundaries. Second,
parallel transport is independent of the charts.


 Except for establishing smoothness of the BV-equivalence maps, 
we have arrived at our result. The flat line bundle that we tensor
with $E$ has the holonomy given by the phase holonomy constructed from
the BV-equivalences as done above. 


We now  establish some smoothness results. We need this to
demonstrate that the transition maps are smooth where they are
supposed to be. This will also establish that the appropriate wave functions in
the domain of \ensuremath{H_{\ensuremath{\kappa}}}, those that are smooth enough to provide a Bohmian
evolution, actually become the smooth sections of $G$ in the
construction. Although this seems like a side remark, it is actually
quite important as it demonstrates the restrictions placed on the
BV-equivalence maps due to property 2 of the definition for being a
BV-equivalent map. 

Given an operator $A$,  let $\ensuremath{\mathcal{D}}(A)$ be the domain of $A$. Define
$\ensuremath{C^{\infty}}(A) \ensuremath{:=} \bigcap_{n=1}^{\infty} \ensuremath{\mathcal{D}}(A^n)$. Let \ensuremath{H_{\ensuremath{\kappa}}}\ be a
cut-independent self-adjoint extension and \ensuremath{H_{\ensuremath{\kappa'}}}\ be a BV-equivalent
extension. Then, from the definition of BV-equivalence, it is clear that
$\ensuremath{C^{\infty}}(\ensuremath{H_{\ensuremath{\kappa}}})$ is mapped into $\ensuremath{C^{\infty}}(\ensuremath{H_{\ensuremath{\kappa'}}})$. By Sobolev regularity
theorems for the Laplacian, we actually have that $\ensuremath{C^{\infty}}(\ensuremath{H_{\ensuremath{\kappa}}})$
consists of smooth functions away from the boundary and these
functions, and their derivatives, have limits as one approaches the
boundaries. Therefore, the image of a smooth section under a 
BV-equivalence map  is smooth. 
Since a BV-equivalence map is multiplication by a
function, we can deduce that the
BV-equivalence map must be smooth away from the boundaries. Furthermore,
the jumps along a given cut are prescribed by the domain of the extension 
corresponding to that cut. 

Since BV-equivalence is an equivalence relation, the same statement
holds for the composition of BV-equivalence maps. Thus, the map
$\ensuremath{U^{\ensuremath{\kappa}{}  \ensuremath{\kappa}{'}}} \ensuremath{U^{\ensuremath{\kappa}{''}  \ensuremath{\kappa}{}}}$ is smooth along $\ensuremath{\kappa}$ wherever $\ensuremath{\kappa}$ does
not intersect either of the other cuts. One can also then understand
that under the equivalence map defining the bundles, sections in
$\ensuremath{C^{\infty}}(\ensuremath{H_{\ensuremath{\kappa}}})$ that
might not look smooth along the boundary do get mapped to sections
that are smooth along \ensuremath{\kappa}\ in the new chart although they will
generally cease to be smooth along the cut of the new chart. 

\end{proof}


\section{Comparisons} {\label{s:comp}}

In this chapter, we have described three very different approaches to
finding various theories. To summarize, we have shown that there is
only one Abelian Quantization Principle, i.e., 
\[
(\ensuremath{\mathcal{Q}}, E, V, \ensuremath{\gamma}) 
\ensuremath{:=} (\ensuremath{\mathcal{Q}}, E, V, \gamma)_{\ensuremath{\widehat{\ensuremath{\mathcal{Q}}}}} 
\ensuremath{\simeq} (\ensuremath{\mathcal{Q}}, E, V, \gamma)_{\nabla} 
\ensuremath{\simeq} (\ensuremath{\mathcal{Q}}, E, V, \gamma)_{H}.
\] 
Although equivalent,  each approach had a different feeling to
it. Indeed, the covering space approach is to change the domain of
the wave function\ into a simply connected space. The bundle approach is to change the
notion of derivative; to a certain extent, we can view this as
changing the range of the wave function. 
The self-adjoint extension
approach directly  
changes the Hamiltonian. It is not at all obvious that all
three approaches should have worked and it is quite remarkable 
 that they produce the same
possibilities. We shall now attempt to contrast the three
approaches. 


Our first approach was to use the covering space; this 
 is a topological solution. Indeed, since 
topology is responsible for these extra possibilities, it seems quite
natural to attempt to trivialize the topology. The usual
topological mechanism for  a trivialization is the use of the 
 simply connected covering space particularly when gradients are involved. 
 We used Bohmian mechanics to explain why the
covering space is useful. The Bohmian constraint is the condition of
projectability of the velocity field. The immediate solution
 is the scalar periodicity condition. We demonstrated that this works
 for any Bohmian theory. In the next chapter, we shall show that,
 generically, the scalar periodicity 
 condition is the only periodicity condition that works.
The true advantage of the covering space is its simplicity. Noticing
 that it is a character which is at work is immediate. It also 
provides a single framework for all of the various
 possible theories.
 A minor disadvantage is that the covering space causes certain
 worries to occur. One worry, easily dismissable, but very
 tempting, is to be bothered by the possibility of 
  wave functions  on the covering space 
 whose Bohmian velocities project down to the base
 space for all time without satisfying any periodicity conditions.
The appropriate response is that the goal is to find theories, not
 exceptions, but many may be unsatisfied with this. Another worry, of
 a very technical nature, is that for spaces with infinite fundamental
 groups, the wave functions satisfying a periodicity condition are not in
 \ensuremath{L^{2}}. This technical problem can be dealt with in a variety of ways,
 such as appealing to the other approaches or perhaps to the theory of
 rigged Hilbert spaces, 
 but it is still an annoyance. 


The approach of considering different self-adjoint extensions is a
functional analysis approach. The issue is, after all, about the
evolution of a system. But in order to consider these possibilities,
we had to the cut space in such a way to make it simply
connected. Then our self-adjoint extension could be thought of as
describing what happens along this virtual boundary. The picture to
have is that 
as a wave packet passes through the boundary, it picks up a phase. But
considering all self-adjoint extensions is to consider a great number
of possible theories. The necessary demand for the extension to be
physically reasonable was that the extension was cut-independent. This
notion was grounded in the Bohmian evolution not detecting which cut
was used. 
 We immediately constructed a cut-independent extension for each
character; one way was to appeal to the covering space. But in this
context, it seemed necessary to rule out the vast number of other
possibilities, which we did.  Unlike the covering space, these
possibilities  did form their own theory. We dismissed them 
on the grounds that  
 they  introduced an element assumed to be not  in the physics, in
much the same way that physics demands an appropriate 
invariance under coordinate changes. 

The bundle approach is a geometric solution. As wave functions are sections of
a vector bundle,  the idea is to consider various different
connections potentially involving various bundles. 
The important fact is that having a non-trivial fundamental group 
implies the  existence of inequivalent flat bundles. 
 This allowed us to define different Bohmian
theories by tensoring the original bundle with a flat line bundle. 
Interestingly, 
the bundle approach does not seem to  offer any wrong directions. 
We had no isolated
wave functions and we had no obvious alternative theories. In terms of
tensoring with a line bundle, only the flat ones are allowed and they
are easily classified. One could ask the question if it is possible to
have locally equivalent bundles which are not constructed in this
manner. This is the case for flat bundles, 
as all flat bundles of a given dimension are
locally equivalent to each other, 
but from the experience with spin as detailed in the next chapter, 
it seems that the only generically
permitted construction is tensoring with a flat line bundle. 


One strength of the bundle approach is in the following question:
 Given the existence
of certain structures, such as spin,
 necessary for the kinds of theories we want,
what are the possible bundles over the configuration space that can
support these structures? 
Although interesting, this question is not what this work is
 addressing.  Indeed, that question applies for simply
connected spaces and requires the consideration of very specific theories
such as those involving spin. 


As to the role of Bohmian mechanics in these considerations, the main
role is to solidify the notion that wave functions
are not primary objects and to explain what is of primary concern,
namely the particles and their evolution.
Indeed, the wave function\ does not need to be well-defined on the base
space, only the Bohmian velocity field needs to be. 
To put it more
succinctly, wave functions should be thought of as more in the nature of
potentials, which are certain artifacts of describing the dynamics,
and less in the nature of electromagnetic fields, which are actually 
exact objects. 

And yet we are not implying that these arguments cannot be made in
standard quantum mechanics. If one accepts, as an axiom, that the local
density is the central object in the theory, then the same results
will follow in almost the same way. 
The cut 
 approach would replace the preservation of the velocity fields
with the preservation of the density. 
In the covering space approach, the projectability of the density would
 replace the projectability of the velocity field. 
The only actual change in the arguments would be the need to use the
Laplacian to demonstrate local constancy of the periodicity or boundary
conditions. 

\section{Local equivalence} \label{is:leqv}

A hidden principle that we have found to be a useful guide is that of
local equivalence. Given the data $(\ensuremath{\mathcal{Q}}, E, V)$, the various 
theories generated by
the Abelian Quantization Principle\ are all locally equivalent to each other. This means that
only global considerations
 can distinguish the theories. 


What does local equivalence 
mean? Roughly, the equations of motion look locally the same. The
objects generating the theory such as $g$, $\nabla$, and $V$ are all
local objects and, in a certain sense,  locally they remain unchanged. 
The global difference comes in from the
character. Another way of saying what local equivalence is, is to
consider coordinates in a neighborhood 
and write down the equations, restricted to that neighborhood, in those
coordinates. The theories are locally equivalent if, in the
appropriate coordinates, the equations are the same. 
 Appropriate coordinates means the following. Pick a
point $\ensuremath{q} \in \ensuremath{\mathcal{Q}}$ and a simply connected neighborhood of
\ensuremath{q}. 
Choose
a basis of the tangent space at $\ensuremath{q}$ and a basis of $E$ at
\ensuremath{q}. Then pick normal coordinates for \ensuremath{\mathcal{Q}}\ 
and parallel transport the basis
of the vector bundle along the geodesics in \ensuremath{\mathcal{Q}}. One can then show
that the equations are the same. 

Local equivalence helps us limit the possibilities. This chapter has
shown that, as far as we can imagine, AQP describes all the possible
locally equivalent theories. But if we allow  the dimension of the
bundle fiber to change, then we can find other theories for which 
the theories are locally equivalent in terms of individual 
components, but they are not,
strictly speaking, 
locally equivalent as the number of components is different. 
More to the point, if $G$ is the  tensor product of  $E$ with some flat bundle
$F$, then there is no local isomorphism between $G$ and $E$ if $F$ is
not a line bundle. Nevertheless, if the initial wave function\ is a product
state, then its evolution will locally look like it is a section of
$E$. The relevance of such possibilities is rather unclear. 


\section{Quantization of classical systems} \label{is:quant}

We finish the chapter by indicating how to quantize a classical system
from a Bohmian point of view. The procedure is little more than Schr\"odinger\
quantization, i.e.\ writing down  Schr\"odinger's equation. One of the uses of Bohmian
mechanics is that this is all that needs to be chosen; the Bohmian
velocity law is always the same. Bohmian mechanics also provides the
appropriate framework in order to recover the classical system in the
classical limit of the quantum theory. Finally, we use the ideas of
this chapter to quantize a classical system whose force is only locally
a gradient, i.e.\ the potential is only well-defined on the universal
covering space.  

\subsection{Newtonian dynamics on a manifold}
We start with a Riemannian configuration space \ensuremath{\mathcal{Q}}, 
with metric $g$, and a
real-valued potential function $V$ defined on \ensuremath{\mathcal{Q}}.
This is
enough information to define a classical evolution on 
\ensuremath{\mathcal{Q}}. Indeed, the Levi-Civita connection, defined by the metric,
allows the acceleration vector to be defined. The metric also provides
a canonical isomorphism between 
 1-forms and tangent vectors. Thus, given $V$, we have
Newton's equation: $\ddot{Q} = - \nabla V$. Equivalently, Newton's 
 equation may be written, in terms of 1-forms, as 
$g(\ddot{Q}, \cdot) = -dV(\cdot)$. An immediate
consequence of this equation is conservation of energy, 
$ \frac{1}{2}g(\dot{Q},\dot{Q}) + V(Q) = E$. 

\subsection{Bohmian Quantization}
The question we address is what the corresponding quantum system
is. The general prescription is that our wave functions are to be
complex-valued functions and the Hamiltonian is 
$-\ensuremath{\Delta}+V$. The wave function\ evolves according to Schr\"odinger's equation\ \eqref{ie:rse}. 
This is Schr\"odinger\ quantization. To complete the theory, we take Bohm's
equation \eqref{ie:rbe} as defining the evolution of the particles. 

This should be contrasted with standard quantum mechanics in which the
theory is only fully specified after the relevant observables are
chosen. 
This can
be tricky since generally the momentum operator does not exist,
e.g.\ the half-line, or the position operator does not exist, e.g.\ the
circle. The relevant data for quantum theories does exist, as can be
deduced from Bohmian mechanics, but the story becomes much more
complicated and requires the finesse of a master in order to do it
correctly without the aid of Bohmian mechanics. 

\subsection{Classical limit}
When one quantizes a system, the main purpose is to form a quantum
theory which, under certain suitable conditions, will behave like the
given system. 
From a
Bohmian perspective, this is formally clear. Indeed, just as in \ensuremath{\mathbb{R}^{n}},
we can write Schr\"odinger's equation, using the polar form $\ensuremath{\psi} = Re^{iS/\ensuremath{\hbar}}$, as two
coupled real equations. One of the equations is the
Hamilton-Jacobi equation for the above classical system with an
additional term, the quantum potential
$(\frac{-\ensuremath{\hbar}^2\ensuremath{\Delta} R}{2R})$. Thus, in any region in which
 this potential may be
considered negligible, the system will behave classically in that
region. Giving a
precise story of when it is negligible is subtle; see \cite{ADGZ02}. 
The main lesson to learn is that in Bohmian mechanics 
the classical limit only involves
demonstrating
 that the trajectories of the actual particles are approximately classical. In
contrast, standard quantum mechanics must explain what objects are
supposed to behave in a way in which the classical world will emerge. 

We conclude this subsection with what this chapter tells us about 
quantization. Although the above quantization procedure does define a
unique theory given $V$, it is not necessarily 
the only Bohmian  theory with that classical limit. 
If $\ensuremath{\mathcal{Q}}$ is multiply connected, then the Abelian Quantization Principle\ tells us that for
each character of the fundamental group, 
we have a different Bohmian theory involving the
potential $V$. Section {\ref{is:leqv}}\ explained that  these theories
are locally equivalent. Since the classical limit is largely based on
a local viewpoint, we can expect that these various theories all have
the same classical limit. To distinguish between them, one must
consider the quantum effects. 

\subsection{Potentials on the covering space} \label{s:pers}


We have been assuming that we have a
time-independent formal Hamiltonian well-defined on \ensuremath{\mathcal{Q}}. One can
ask the question about quantizing systems with a potential which is
only defined on the covering space; an example is having a constant
electric field along a circle. Alternatively, we could also
 pursue time-dependent, topologically motivated 
connections, such as having a time-dependent magnetic field in the
Aharonov-Bohm setup. On the surface, these seem to be areas where the
approaches are inequivalent. On closer inspection, these are, in fact,
the same phenomena. 


We start with \ensuremath{\mathcal{Q}}\ and a 1-form $\omega$ which is closed but not
exact on \ensuremath{\mathcal{Q}}. 
The model spaces to think about are 
the circle and the torus. 

We take $\omega$ to be our classical force, and we
have its potential $V$ on the covering space, but not on the base
space.  We therefore consider the Bohmian system on the covering space
which corresponds to $-\ensuremath{\Delta} +V$, and the relevant set of wave functions will be
those that consistently project down. We claim that the periodicity
condition now becomes time-dependent. More to the point, our space of
wave functions will be a disjoint union over the periodicity conditions and any
particular wave function\ will evolve from one sector to the next. Note that the
space of  wave functions  ceases to be a linear space. 

To establish what we just claimed, we start with a wave function\ on the
covering space satisfying 
\begin{equation} \label{ei:tpercon}
\ensuremath{\psi} (\ensuremath{\sigma} \ensuremath{\hat{q}}, 0) = \ensuremath{\gamma_{\ensuremath{\sigma}}} \ensuremath{\psi}(\ensuremath{\hat{q}},0).
\end{equation}
 Our claim is that if $V(\ensuremath{\sigma} \ensuremath{\hat{q}}) = V(\ensuremath{\hat{q}}) +
\beta_{\ensuremath{\sigma}}$ for all \ensuremath{\hat{q}}\ and \ensuremath{\sigma}, then 
\begin{equation} \label{ei:tcons}
\ensuremath{\psi} (\ensuremath{\sigma} \ensuremath{\hat{q}} , t) 
= \ensuremath{\gamma_{\ensuremath{\sigma}}} e^{-i \beta_{\ensuremath{\sigma}}t} \ensuremath{\psi} (\ensuremath{\hat{q}}, t).    
\end{equation}
We can establish this by noting that if $\phi$ is any solution to 
Schr\"odinger's equation, then  
\[\phi_{\ensuremath{\sigma}} (\ensuremath{\hat{q}},t)
\ensuremath{:=} 
\ensuremath{\gamma_{\ensuremath{\sigma}}}^{-1}  e^{i \beta_{\ensuremath{\sigma}}t} \phi (\ensuremath{\sigma}
\ensuremath{\hat{q}}, t)
\] is also a solution to Schr\"odinger's equation. That is a simple computation. If \ensuremath{\psi}\
satisfies equation \eqref{ei:tpercon}, then $\ensuremath{\psi}_{\ensuremath{\sigma}}(\ensuremath{\hat{q}}, 0) =
 \ensuremath{\psi}(\ensuremath{\hat{q}}, 0)$. Hence, by uniqueness of solutions, 
\[\ensuremath{\psi}_{\sigma}(\ensuremath{\hat{q}}, t) = \ensuremath{\psi}(\ensuremath{\hat{q}}, t)\] which is equation
 \eqref{ei:tcons}.
That is to say,  we have time-dependent periodicity
conditions. Furthermore, as above, in the classical
limit the motion on \ensuremath{\mathcal{Q}}\ should look Newtonian with force $\omega$.

We now wish to argue that the above Bohmian system
 is equivalent to a Hamiltonian with a
time-dependent connection, but with a fixed periodicity
condition. Quite generally, we claim that the 
Bohmian
systems 
\begin{equation} \label{e:2thry}
(-\ensuremath{\Delta} + V, \nabla) \ensuremath{\sim} (-(\nabla -i
tdV )^2, \nabla -i
tdV )
\end{equation}
are equivalent
under the map $\ensuremath{\psi} \mapsto e^{ i V t} \ensuremath{\psi}$.\footnote{If $V$ depends on
$t$, then we would replace $Vt$ with $W \ensuremath{:=} \int V dt$ and replace
$tdV $ with $dW$. 
}
 Again, a simple computation
establishes this. If we make this change, then the time-dependent
mapping cancels the time-dependent change of phase in the periodicity
condition. Hence we are  left with the initial periodicity condition as
the condition for all $t$. We may 
thus
define a flat bundle over \ensuremath{\mathcal{Q}}, by using the time-independent periodicity
condition. As $dV$ is well-defined on \ensuremath{\mathcal{Q}},  
the time-dependent connection on the bundle over the covering
space induces a time-dependent connection on the bundle over
\ensuremath{\mathcal{Q}}. 

As for the cuts, it seems that time-dependent boundary conditions
would be necessary. Our scheme would be the following. Cut the space
and define $V$ up to the boundary.
If we use time-independent conditions,
then the potential we are modeling is the one corresponding to the
lift of this discontinuous potential to the covering space. The cut
would certainly matter and  presumably no 
BV-equivalence exists between the different cuts. 
However, if we use a time-dependent
boundary condition, as suggested by the time-dependent periodicity condition,
then the extension will be cut-independent.


Let us say a few words about this in relation to a quantum system on
the circle. 
On the circle, there are closed, but not
exact, 1-forms. 
Indeed, a constant force along
the circle is represented by the the non-exact 1-form
$\frac{\beta}{2\pi}d\theta $. We shall suppress the $d\theta$ as is
customarily done.   
On the covering space we
have a potential function, $\frac{\beta}{2\pi} \theta$,  and, therefore, a
well-defined Schr\"odinger's equation.  The
  periodicity condition is $\ensuremath{\psi} (\theta +2\pi n, t) = \gamma^n e^{i
\beta n t} \ensuremath{\psi}(\theta, t)$. 
Such wave functions obviously define a Bohmian evolution on $S^1$.
The change of connection which eliminates this time-dependent
periodicity condition is $\ensuremath{\frac{\partial {}}{\partial {\theta}}} - i \frac{\beta}{2 \pi}
t$. Forming the induced line bundle over the circle, we have the
trivial bundle with 
 connection 
\begin{equation*} 
\nabla \ensuremath{:=} \ensuremath{\frac{\partial {}}{\partial {\theta}}} - i ( \frac{\alpha + t \beta}{2 \pi} )
\end{equation*}
 where $\alpha$
satisfies $\gamma = e^{i \alpha}$. This connection incorporates both the
periodicity condition and the time-dependent connection on the
covering space.

\clearpage
\bibliographystyle{abbrv}
\bibliography{references}
\end{document}
